\subsection{Saturated classes}
\label{subsec:2dloc-saturation}

\gutentag{005A}%
In the setting of \cref{subsec:2dloc-bousfield}, the class of $S$-local equivalences (the $f$ such that every $S$-local object is $\{f\}$-local) is strongly saturated: it contains the equivalences, is closed under pushouts and under colimits in $\func(\Delta^{1},\C)$, and satisfies two-out-of-three. For presentable $\C$ and a set $S$ it is the strongly saturated class generated by $S$ \cite[5.5.4.15]{htt}, so that accessible localizations and strongly saturated classes of small generation determine each other.
We will categorify the first half in an arbitrary $(\infty,2)$-category $\D$, and reduce the second half to a statement about the units of the localization (\cref{thm:sat-units}). Three things change.
\begin{itemize}
\gutentag{005B}%
\item Two-out-of-three will be replaced by left cancellation: if $g$ and $gf$ are in the class, so is $f$. Right cancellation fails (\cref{ex:saturation-fails}).
\item Colimits will be allowed only along Beck--Chevalley squares (\cref{ex:saturation-fails}). So a saturated class will consist of $1$-morphisms together with commutative squares between them, i.e.\ it will be a locally full sub-$(\infty,2)$-category of $\func(\Delta^{1},\D)$.
\item There are two notions, one for coreflections and one for reflections, exchanged by $(-)^{\mathrm{co}}$. They have the same closure axioms and differ only in the base case: coreflective $1$-morphisms with right Beck--Chevalley squares, resp.\ reflective ones with left Beck--Chevalley squares (\cref{lem:mates-refl}).
\end{itemize}
\gutentag{005C}%
Every locally full sub-$(\infty,2)$-category $\C \subset \D$ will have a coreflective and a reflective saturation $\Sat^{\mathrm{corefl}}_{\C}$, $\Sat^{\mathrm{refl}}_{\C}$ (\cref{def:saturation-C}), and the saturation $\overline{P}$ of a class $P$ (\cref{def:saturation}) will be the class of $1$-morphisms of $\Sat^{\mathrm{corefl}}_{\C}$ for $\C = \Pc_{\emptyset,P}$.

\subsubsection*{Coreflective $1$-morphisms}

\gutentag{005D}%
A $1$-morphism $g\colon y \to z$ is \emph{fully faithful} if every $g_{\natural}\colon \Hom(t,y) \to \Hom(t,z)$ is. Reflective and coreflective $1$-morphisms are fully faithful (\cref{def:reflection}).

\begin{lem}[cancellation]
\label{lem:ff-cancel}
\gutentag{005E}%
Let $f\colon x \to y$ and let $g\colon y \to z$ be fully faithful. If $gf \dashv r$, then $f \dashv rg$, with the unit of $gf \dashv r$. Dually, if $l \dashv gf$, then $lg \dashv f$, with the counit of $l \dashv gf$.
In particular, if $g$ and $gf$ are coreflective, so is $f$; and if $g$ and $gf$ are reflective, so is $f$.
\end{lem}

\gutentag{005F}%
\reviewcomment{When we write an introduction, refer to here as the reason why we want (co)reflections
rather than adjoints.}

\begin{proof}
\gutentag{0060}%
Let $\eta$, $\varepsilon$ be the unit and counit of $gf \dashv r$. As $g_{\natural}$ is fully faithful on $\Hom(y,y)$, there is a unique $\varepsilon'\colon frg \Rightarrow \mathrm{id}_{y}$ with $g\varepsilon' = \varepsilon g$. The triangle identities for $(\eta,\varepsilon')$ hold in the homotopy $2$-category, as $g_{\natural}$ is faithful and
\[
g(\varepsilon'f \cdot f\eta) = \varepsilon gf \cdot gf\eta = \mathrm{id}_{gf} \ , \qquad rg\varepsilon' \cdot \eta rg = (r\varepsilon \cdot \eta r)g = \mathrm{id}_{rg} \ ,
\]
so $\eta$ and $\varepsilon'$ are the unit and counit of an adjunction $f \dashv rg$ \cite[Thm.~4.3.11]{riehlverity2016adjunctions}. The dual follows by applying $(-)^{\mathrm{co}}$.
\end{proof}

\begin{lem}
\label{lem:corefl-pullback}
\gutentag{0061}%
A functor of $\infty$-categories $R\colon B \to A$ has a fully faithful left adjoint if and only if its pullback $C \times_{A} B \to C$ along every $F\colon C \to A$ is a right adjoint. In that case, if $\Lambda \dashv R$, the pullback has the fully faithful left adjoint $c \mapsto (c,\Lambda Fc)$.
\end{lem}

\begin{proof}
\gutentag{0062}%
If the unit of $\Lambda \dashv R$ is invertible, then
\[
\mapp\bigl((c,\Lambda Fc),(c',b')\bigr) \simeq \mapp(c,c') \times_{\mapp(Fc,Rb')} \mapp(\Lambda Fc,b') \simeq \mapp(c,c') \ ,
\]
as $R\colon \mapp(\Lambda Fc,b') \to \mapp(R\Lambda Fc,Rb') \simeq \mapp(Fc,Rb')$ is an equivalence.
Conversely, pulling back along $a\colon \mathrm{pt} \to A$ shows that the fibre $B_{a}$ has an initial object $b_{a}$. For $\varphi\colon a \to Rb'$, let $B' := \Delta^{1} \times_{A} B$ be the pullback along $\varphi$. The left adjoint of $B' \to \Delta^{1}$ sends $0$ to an initial object of $B'$, which maps to $b_{a} \in B'_{0}$ and hence lies in $B'_{0} = B_{a}$; so $b_{a}$ is initial in $B'$. Thus the fibre of $R\colon \mapp_{B}(b_{a},b') \to \mapp_{A}(a,Rb')$ over $\varphi$, which is $\mapp_{B'}(b_{a},b')$, is contractible. So $\mathrm{id}\colon a \simeq Rb_{a}$ is a universal arrow, and $R$ has a left adjoint with invertible unit.
\end{proof}

\gutentag{0063}%
So the right adjoints of coreflections are exactly the functors that remain right adjoints after every base change: coreflections are the base-change-stable version of adjunctions. For $1$-morphisms of $\D$, the corresponding statement will be that coreflective $1$-morphisms are stable under $(\infty,2)$-pushouts (\cref{thm:saturation-saturated} for $\C = \D$), while left adjoints are not (\cref{ex:saturation-fails}).

\subsubsection*{Weighted colimits}
\gutentag{0064}%
For a small $(\infty,2)$-category $K$ and functors $W, G\colon K\opcat \to \Cat$ write $\{W,G\} := \Hom_{\func(K\opcat,\Cat)}(W,G) \in \Cat$, the $W$-weighted limit of $G$.
Let $F\colon K \to \D$. A \emph{$W$-weighted colimit} of $F$ is an object $W \star F$ with a $1$-morphism $\kappa\colon W \Rightarrow \Hom_{\D}(F(-),W \star F)$ in $\func(K\opcat,\Cat)$ such that
\[
\Hom_{\D}(W \star F,e) \xrightarrow{\ \kappa^{\natural} \circ (-)_{\natural}\ } \{W,\Hom_{\D}(F(-),e)\}
\]
is an equivalence for every $e \in \D$. It is natural in $e$.
For $W$ constant with value $\mathrm{pt}$ these are the conical $(\infty,2)$-colimits, e.g.\ the $(\infty,2)$-pushouts of \cref{subsec:2dloc-plain}; for $K = \mathrm{pt}$ and $W = A \in \Cat$ it is the tensor $A \otimes x$; and the cylinder $\Cyl(s)$ of \cref{constr:cylinder} is the colimit of $x \xleftarrow{\mathrm{id}} x \xrightarrow{s} y$ weighted by $\mathrm{pt} \xrightarrow{0} \Delta^{1} \xleftarrow{1} \mathrm{pt}$.

\gutentag{0065}%
A functor $F\colon K \to \func(\Delta^{1},\D)$ is a transformation $F_{0} \Rightarrow F_{1}$ of functors $K \to \D$. If $W \star F_{0}$ and $W \star F_{1}$ exist, it induces $f\colon W \star F_{0} \to W \star F_{1}$, and for $k \in K$ and $w \in W(k)$ the components of $\kappa$ at $w$ form a commutative square $F(k) \to f$, which we call a \emph{colimit square}.

\subsubsection*{Saturated classes}

\begin{defi}
\label{def:saturated}
\gutentag{0066}%
Let $\Sat \subset \func(\Delta^{1},\D)$ be a locally full sub-$(\infty,2)$-category. We say that a $1$-morphism of $\D$ \emph{is in $\Sat$} if it is an object of $\Sat$, and call the $1$-morphisms of $\Sat$, which are commutative squares, \emph{$\Sat$-squares}. $\Sat$ is \emph{coreflectively saturated} if:
\begin{enumerate}
\item[(S1)] every coreflective $1$-morphism is in $\Sat$, and every right Beck--Chevalley square between coreflective $1$-morphisms is an $\Sat$-square; equivalently, $\ell^{*}\colon \func(\Corefl,\D) \to \func(\Delta^{1},\D)$ factors through $\Sat$ (\cref{lem:mates-refl});
\item[(S2)] (composition) if $f\colon x \to y$ and $g\colon y \to z$ are in $\Sat$, then so is $gf$, and the square $(\mathrm{id}_{x},g)\colon f \to gf$ is an $\Sat$-square;
\item[(S3)] (left cancellation) if $g$ and $gf$ are in $\Sat$, then so is $f$;
\item[(S4)] (colimits) for every functor $F\colon K \to \Sat$ from a small $(\infty,2)$-category and every $W\colon K\opcat \to \Cat$ such that $W \star F_{0}$ and $W \star F_{1}$ exist, the induced $f$ is in $\Sat$ and the colimit squares $F(k) \to f$ are $\Sat$-squares.
\end{enumerate}
$\Sat$ is \emph{reflectively saturated} if it satisfies (S2)--(S4) and
\begin{enumerate}
\item[(S1')] every reflective $1$-morphism is in $\Sat$, and every left Beck--Chevalley square between reflective $1$-morphisms is an $\Sat$-square; equivalently, $\rho^{*}\colon \func(\Refl,\D) \to \func(\Delta^{1},\D)$ factors through $\Sat$.
\end{enumerate}
\end{defi}

\gutentag{0067}%
Here the square $(\mathrm{id}_{x},g)$ in (S2) has $\mathrm{id}_{x}$ as its component on sources and $g$ on targets. For $f = \mathrm{id}_{x}$ it is the square $\mathrm{id}_{x} \to g$. The Beck--Chevalley conditions in (S1) and (S1') are those of \cref{lem:mates}(1) and (2).
The conditions are closure conditions, so intersections of (co)reflectively saturated classes are (co)reflectively saturated. For a class $P$ of $1$-morphisms, or more generally a diagram $K \to \func(\Delta^{1},\D)$ of $1$-morphisms and squares, we write $\langle P \rangle$ for the coreflectively saturated class it generates.

\gutentag{0068}%
$(-)^{\mathrm{co}}$ exchanges the two notions: $\func(\Delta^{1},\D^{\mathrm{co}}) \simeq \func(\Delta^{1},\D)^{\mathrm{co}}$ has the same objects and $1$-morphisms as $\func(\Delta^{1},\D)$, and $(-)^{\mathrm{co}}$ exchanges reflective and coreflective $1$-morphisms and left and right Beck--Chevalley squares (\cref{rmk:adjoin-duality}). (S2) and (S3) are self-dual, and so is (S4): as $\Hom_{\D^{\mathrm{co}}}(c,e) = \Hom_{\D}(c,e)\opcat$, a $W$-weighted colimit of $F\colon K \to \D^{\mathrm{co}}$ is a $((-)\opcat \circ W)$-weighted colimit of the corresponding functor $K^{\mathrm{co}} \to \D$.
So reflections and coreflections need separate definitions only in the base case (S1), (S1').

\begin{exmp}
\label{ex:saturation-fails}
\gutentag{0069}%
Let $\D = \Cat$, where coreflective $1$-morphisms are the fully faithful left adjoints; pushouts in $\Cat$ are $(\infty,2)$-pushouts.
\begin{enumerate}
\item (No right cancellation.) $f = \{0\}\colon \mathrm{pt} \to \Delta^{1}$ and $gf = \mathrm{id}_{\mathrm{pt}}$, for $g\colon \Delta^{1} \to \mathrm{pt}$, are coreflective, but $g$ is not fully faithful.
\item (Colimits need Beck--Chevalley squares.) The square $\sigma\colon \{0\} \to \mathrm{id}_{\Delta^{1}}$ with components $\{0\}\colon \mathrm{pt} \to \Delta^{1}$ on sources and $\mathrm{id}_{\Delta^{1}}$ on targets is not right Beck--Chevalley: its right mate is the non-invertible $2$-morphism $\{0\} \circ (\Delta^{1} \to \mathrm{pt}) \Rightarrow \mathrm{id}_{\Delta^{1}}$. The square $\tau\colon \{0\} \to \mathrm{id}_{\mathrm{pt}}$ with components $\mathrm{id}_{\mathrm{pt}}$ and $\Delta^{1} \to \mathrm{pt}$ is right Beck--Chevalley, as its mate is a $2$-morphism between functors to $\mathrm{pt}$. The pushout of $\mathrm{id}_{\Delta^{1}} \xleftarrow{\sigma} \{0\} \xrightarrow{\tau} \mathrm{id}_{\mathrm{pt}}$ in $\func(\Delta^{1},\Cat)$ is
\[
\Delta^{1} \sqcup_{\mathrm{pt}} \mathrm{pt} = \Delta^{1} \longrightarrow \mathrm{pt} = \Delta^{1} \sqcup_{\Delta^{1}} \mathrm{pt} \ ,
\]
which is a left adjoint but not coreflective, although all three $1$-morphisms in the diagram are coreflective.
\item (Plain adjoints.) Let $\overline{\emptyset}^{\dashv}$ be the class of left adjoint functors. Left cancellation fails for $f = \{1\}\colon \mathrm{pt} \to \Delta^{1}$ and $g\colon \Delta^{1} \to \mathrm{pt}$: $g$ and $gf = \mathrm{id}$ are left adjoints, but $f$ is not, as $1$ is not initial. Pushouts: the pushout of $g$ along the inclusion of $c \to c'$ into $\Lambda := \{c' \leftarrow c \to d\}$ is the functor $\Lambda \to \Delta^{1}$ sending $c,c'$ to $0$ and $d$ to $1$, which is not a left adjoint as $\Lambda$ has no terminal object.
\end{enumerate}
\end{exmp}

\subsubsection*{The saturation of a locally full sub-$(\infty,2)$-category}

\begin{defi}
\label{def:saturation}
\gutentag{006A}%
The \emph{saturation} of a class $P$ of $1$-morphisms of $\D$ is the class $\overline{P}$ of $1$-morphisms $f$ such that $\Pc_{\emptyset,P \cup \{f\}} = \Pc_{\emptyset,P}$. Likewise $\overline{P}^{\dashv}$ is the class of $f$ with $\Pc^{\dashv}_{\emptyset,P \cup \{f\}} = \Pc^{\dashv}_{\emptyset,P}$.
\end{defi}

\gutentag{006B}%
By \cref{lem:repr-refl}, $\overline{\emptyset}$ is the class of coreflective $1$-morphisms of $\D$, and $\overline{\emptyset}^{\dashv}$ that of left adjoints. We now describe $\overline{P}$ together with squares.

\gutentag{006C}%
Let $\iota\colon \C \hookrightarrow \D$ be a locally full sub-$(\infty,2)$-category. By the enriched Yoneda lemma \cite{hinich2020yoneda,heine2023enriched},
\[
h_{\C}\colon \D\opcat \longrightarrow \func(\C,\Cat) \ , \qquad d \mapsto \Hom_{\D}(d,\iota-) \ , \quad f \mapsto f^{\natural} \ ,
\]
is a functor. Reversing the arrow $\Delta^{1}$ gives $\func(\Delta^{1},\D)\opcat \simeq \func(\Delta^{1},\D\opcat)$, under which $f\colon x \to y$ corresponds to $f\opcat\colon y \to x$, sent by $h_{\C}$ to $f^{\natural}$.
Recall that the equivalence $\Adj \simeq \Adj\opcat$, $l \mapsto r\opcat$, $r \mapsto l\opcat$, preserves units and counits, so it induces $\Corefl \simeq \Corefl\opcat$ and $\Refl \simeq \Refl\opcat$ carrying $\ell$ to $\rho\opcat$: a $1$-morphism $l$ of $\D$ is coreflective if and only if $l\opcat$ is the right adjoint of a coreflection in $\D\opcat$.

\begin{defi}
\label{def:saturation-C}
\gutentag{006D}%
The \emph{coreflective saturation} $\Sat^{\mathrm{corefl}}_{\C}$ of $\C$ is the locally full sub-$(\infty,2)$-category of $\func(\Delta^{1},\D)$ whose opposite is the pullback of
\[
\rho^{*}\colon \func\bigl(\Corefl,\func(\C,\Cat)\bigr) \longrightarrow \func\bigl(\Delta^{1},\func(\C,\Cat)\bigr)
\]
along $\func(\Delta^{1},\D)\opcat \simeq \func(\Delta^{1},\D\opcat) \xrightarrow{h_{\C}} \func(\Delta^{1},\func(\C,\Cat))$. The \emph{reflective saturation} $\Sat^{\mathrm{refl}}_{\C}$ is defined in the same way with $\ell^{*}\colon \func(\Refl,\func(\C,\Cat)) \to \func(\Delta^{1},\func(\C,\Cat))$.
\end{defi}

\gutentag{006E}%
Both are locally full, since $\rho^{*}$ and $\ell^{*}$ are locally full inclusions (\cref{lem:mates-refl}), these are stable under pullback, and $(-)\opcat$ preserves them.

\begin{lem}
\label{lem:saturation-explicit}
\gutentag{006F}%
\leavevmode
\begin{enumerate}
\item A $1$-morphism $f\colon x \to y$ is in $\Sat^{\mathrm{corefl}}_{\C}$ if and only if every object and $1$-morphism of $\C$ is $(\emptyset,\{f\})$-local (\cref{def:IP-local}): every $f^{\natural}\colon \Hom_{\D}(y,\iota q) \to \Hom_{\D}(x,\iota q)$ has a fully faithful left adjoint $\lambda_{f,q}$, and $g_{\natural}$ commutes with these for every $1$-morphism $g$ of $\C$.
\item Let $f\colon x \to y$ and $f'\colon x' \to y'$ be in $\Sat^{\mathrm{corefl}}_{\C}$. A commutative square $\sigma\colon f \to f'$ with components $a\colon x \to x'$, $b\colon y \to y'$ is an $\Sat^{\mathrm{corefl}}_{\C}$-square if and only if for every $q \in \C$ the square
\[
\begin{tikzcd}[column sep=large]
\Hom_{\D}(y',\iota q) \ar[r, "f'^{\natural}"] \ar[d, "b^{\natural}"'] & \Hom_{\D}(x',\iota q) \ar[d, "a^{\natural}"] \ar[dl, Rightarrow, shorten=4mm] \\
\Hom_{\D}(y,\iota q) \ar[r, "f^{\natural}"'] & \Hom_{\D}(x,\iota q)
\end{tikzcd}
\]
(filled by $\sigma$) is left Beck--Chevalley, i.e.\ $\lambda_{f,q}a^{\natural} \Rightarrow b^{\natural}\lambda_{f',q}$ is invertible: extending along $f$ commutes with restriction along $\sigma$.
\end{enumerate}
Dually, $f$ is in $\Sat^{\mathrm{refl}}_{\C}$ if and only if every object and $1$-morphism of $\C$ is $(\{f\},\emptyset)$-local, and $\Sat^{\mathrm{refl}}_{\C}$-squares are those for which the squares above are right Beck--Chevalley.
In particular, for $\C = \Pc_{I,P}$, $f$ is in $\Sat^{\mathrm{corefl}}_{\C}$ if and only if $\Pc_{I,P \cup \{f\}} = \Pc_{I,P}$, and in $\Sat^{\mathrm{refl}}_{\C}$ if and only if $\Pc_{I \cup \{f\},P} = \Pc_{I,P}$. For $I = \emptyset$ the $1$-morphisms in $\Sat^{\mathrm{corefl}}_{\C}$ form the class $\overline{P}$ of \cref{def:saturation}.
\end{lem}

\begin{proof}
\gutentag{0070}%
By \cref{lem:mates-refl}(2), $f$ is in $\Sat^{\mathrm{corefl}}_{\C}$ if and only if $f^{\natural}$ is the right adjoint of a coreflection in $\func(\C,\Cat)$. By \cite[Cor.~5.2.12]{abellangagnahaugseng2024straightening}, $f^{\natural}$ is a right adjoint if and only if its components are and its transposed naturality squares, which are the second squares of \cref{def:IP-local}, are left Beck--Chevalley; its unit is invertible if and only if its components are. This is (1).
By \cref{lem:mates-refl}(2) and \cref{lem:mates}(2), $\sigma$ is an $\Sat^{\mathrm{corefl}}_{\C}$-square if and only if the square of the statement, viewed in $\func(\C,\Cat)$, is left Beck--Chevalley. Evaluation at $q$ preserves adjunctions and hence mates, and a $2$-morphism of $\func(\C,\Cat)$ is invertible if and only if its components are. This is (2). The reflective case is the same with \cref{lem:mates}(1).
The last claims follow from (1), as $\Pc_{I,P \cup \{f\}} = \Pc_{I,P} \cap \Pc_{\emptyset,\{f\}}$.
\end{proof}

\begin{lem}
\label{lem:ff-reflects}
\gutentag{0071}%
Let $G\colon \A \to \B$ be a fully faithful functor, i.e.\ every $\Hom_{\A}(a,a') \to \Hom_{\B}(Ga,Ga')$ is an equivalence. For $X \in \{\Refl,\Corefl\}$ the square
\[
\begin{tikzcd}
\func(X,\A) \ar[r] \ar[d, "\ell^{*}"'] & \func(X,\B) \ar[d, "\ell^{*}"] \\
\func(\Delta^{1},\A) \ar[r] & \func(\Delta^{1},\B)
\end{tikzcd}
\]
is a pullback, and likewise with $\rho^{*}$: $G$ reflects (co)reflective $1$-morphisms and Beck--Chevalley squares between them.
\end{lem}

\begin{proof}
\gutentag{0072}%
Both vertical maps are locally full inclusions (\cref{lem:mates-refl}), so it suffices to compare objects and $1$-morphisms. Being a left or right adjoint, and being reflective or coreflective, are detected in the homotopy $2$-category \cite[Thm.~4.3.11]{riehlverity2016adjunctions}. As $G$ is an equivalence on hom-categories, an adjunction in $\htwo\B$ between $1$-morphisms $Ga \to Ga'$, $Ga' \to Ga$ is the image of one in $\htwo\A$, with invertible unit or counit if and only if its image has one. $G$ carries mates to mates and reflects invertibility of $2$-morphisms.
\end{proof}

\begin{prop}
\label{prop:saturation-L}
\gutentag{0073}%
\leavevmode
\begin{enumerate}
\item $\Sat^{\mathrm{corefl}}_{\D}$ is the image of $\ell^{*}\colon \func(\Corefl,\D) \to \func(\Delta^{1},\D)$: the coreflective $1$-morphisms and the right Beck--Chevalley squares between them. Likewise $\Sat^{\mathrm{refl}}_{\D}$ is the image of $\rho^{*}\colon \func(\Refl,\D) \to \func(\Delta^{1},\D)$.
\item Suppose $\iota\colon \C \to \D$ has a left adjoint $L$ (\cref{def:enriched-adjoint}). Then $\Sat^{\mathrm{corefl}}_{\C}$ is the preimage of the image of $\ell^{*}\colon \func(\Corefl,\C) \to \func(\Delta^{1},\C)$ under $L$: $f$ is in $\Sat^{\mathrm{corefl}}_{\C}$ if and only if $Lf$ is coreflective, and a square $\sigma$ between such $1$-morphisms is an $\Sat^{\mathrm{corefl}}_{\C}$-square if and only if $L\sigma$ is right Beck--Chevalley. Likewise for $\Sat^{\mathrm{refl}}_{\C}$ with reflective $1$-morphisms and left Beck--Chevalley squares.
\end{enumerate}
\end{prop}

\begin{proof}
\gutentag{0074}%
(1) $h_{\D}$ is the enriched Yoneda embedding, which is fully faithful \cite{hinich2020yoneda,heine2023enriched}. By \cref{lem:ff-reflects} for $X = \Corefl$ and $\rho^{*}$, a $1$-morphism or square of $\func(\Delta^{1},\D\opcat)$ lies in the image of $\rho^{*}\colon \func(\Corefl,\D\opcat) \to \func(\Delta^{1},\D\opcat)$ if and only if its image under $h_{\D}$ does. Under $\Corefl \simeq \Corefl\opcat$ the former image is the opposite of the image of $\ell^{*}\colon \func(\Corefl,\D) \to \func(\Delta^{1},\D)$.
(2) The equivalences $\Hom_{\C}(Ld,q) \simeq \Hom_{\D}(d,\iota q)$ are natural in $d$ and $q$, so $h_{\C} \simeq h'_{\C} \circ L\opcat$ for the enriched Yoneda embedding $h'_{\C}\colon \C\opcat \to \func(\C,\Cat)$. So $\Sat^{\mathrm{corefl}}_{\C}$ is the preimage under $L$ of the coreflective saturation of $\C$ in itself, which is the image of $\ell^{*}$ by (1) for $\C$.
\end{proof}

\gutentag{0075}%
For presentable $\D$ and a set $P$, (2) says that $\overline{P}$ consists of the $f$ such that $Lf$ is coreflective in $\D\corefl{P}$: the analogue of the $L$-equivalences. We will record this in \cref{cor:presentable}.

\begin{thm}
\label{thm:saturation-saturated}
\gutentag{0076}%
For every locally full sub-$(\infty,2)$-category $\C \subset \D$, $\Sat^{\mathrm{corefl}}_{\C}$ is coreflectively saturated and $\Sat^{\mathrm{refl}}_{\C}$ is reflectively saturated.
In particular (for $\C = \D$, by \cref{prop:saturation-L}(1)) the coreflective $1$-morphisms with the right Beck--Chevalley squares form the smallest coreflectively saturated class, and the reflective $1$-morphisms with the left Beck--Chevalley squares the smallest reflectively saturated class.
\end{thm}

\begin{proof}
\gutentag{0077}%
We treat $\Sat := \Sat^{\mathrm{corefl}}_{\C}$; the reflective case is the same with $\Refl$ and $\ell^{*}$ in place of $\Corefl$ and $\rho^{*}$. Write $\E := \func(\C,\Cat)$ and $h := h_{\C}$.

(S1) Under $\Corefl \simeq \Corefl\opcat$, the image of $\ell^{*}\colon \func(\Corefl,\D) \to \func(\Delta^{1},\D)$ is the opposite of the image of $\rho^{*}\colon \func(\Corefl,\D\opcat) \to \func(\Delta^{1},\D\opcat)$. As $h$ is a functor, it commutes with $\rho^{*}$ and so carries the latter into the image of $\rho^{*}$ in $\func(\Delta^{1},\E)$.

(S2) Let $\lambda_{f} \dashv f^{\natural}$ and $\lambda_{g} \dashv g^{\natural}$ in $\E$ have invertible units. Then $\lambda_{g}\lambda_{f} \dashv f^{\natural}g^{\natural} = (gf)^{\natural}$ has invertible unit, so $gf$ is in $\Sat$. The square $(\mathrm{id}_{x},g)$ is sent to the square with $u = (gf)^{\natural}$, $u' = f^{\natural}$, $f_{0} = g^{\natural}$, $f_{1} = \mathrm{id}$ in the notation of \cref{lem:mates}(2), filled by the identity. By the triangle identities its left mate $\lambda_{f} \Rightarrow g^{\natural}\lambda_{g}\lambda_{f}$ is the unit of $\lambda_{g} \dashv g^{\natural}$ whiskered with $\lambda_{f}$, which is invertible.

(S3) In $\E\opcat$ the right adjoints of coreflections become coreflective $1$-morphisms (as above), and $(gf)^{\natural} = f^{\natural}g^{\natural}$ becomes the composite of $(f^{\natural})\opcat$ followed by $(g^{\natural})\opcat$. So \cref{lem:ff-cancel} in $\E\opcat$, with $(g^{\natural})\opcat$ coreflective and hence fully faithful, shows that $(f^{\natural})\opcat$ is coreflective.

(S4) Let $F\colon K \to \Sat$ and $W$ be as in (S4), and let $x := W \star F_{0}$, $y := W \star F_{1}$. By the definition of weighted colimits, naturally in $q \in \C$,
\[
\Hom_{\D}(x,\iota q) \simeq \{W,\Hom_{\D}(F_{0}(-),\iota q)\} \ , \qquad \Hom_{\D}(y,\iota q) \simeq \{W,\Hom_{\D}(F_{1}(-),\iota q)\} \ ,
\]
compatibly with $f^{\natural}$ and with the colimit squares. So, writing $H\colon K\opcat \to \func(\Delta^{1},\E) \simeq \func(\Delta^{1} \times \C,\Cat)$ for $h \circ F\opcat$, we have $f^{\natural} \simeq \{W,H\}$, formed pointwise in $\Delta^{1} \times \C$, and the colimit square at $w \in W(k)$ is sent to the evaluation $\{W,H\} \to H(k)$ at $w$.
As $F$ lands in $\Sat$, $H$ sends objects and $1$-morphisms of $K\opcat$ into the image of the locally full inclusion $\rho^{*}$, so it lifts to $G\colon K\opcat \to \func(\Corefl,\E) \simeq \func(\Corefl \times \C,\Cat)$ \cite[Cor.~2.6.6]{abellangagnahaugseng2024straightening}. Let $\{W,G\} \in \func(\Corefl \times \C,\Cat)$ be formed pointwise, i.e.\ the composite of $\Corefl \times \C \to \func(K\opcat,\Cat)$, adjoint to $G$, with $\{W,-\}$. As $\rho^{*}$ is restriction, $\rho^{*}\{W,G\} = \{W,H\} \simeq f^{\natural}$, and $\rho^{*}$ carries the evaluations $\{W,G\} \to G(k)$ at $w \in W(k)$, which are $1$-morphisms of $\func(\Corefl,\E)$, to the images of the colimit squares. So $f$ is in $\Sat$ and the colimit squares are $\Sat$-squares.
\end{proof}

\gutentag{0078}%
Taking $K = \Lambda^{2}_{0}$, $W = \mathrm{pt}$, and the diagram $f \leftarrow \mathrm{id}_{a} \rightarrow \mathrm{id}_{c}$ given by the $\Sat$-squares $(\mathrm{id}_{a},f)$ of (S2) and $(g,g)$ of (S1), (S4) says that $\Sat$ is stable under $(\infty,2)$-pushouts along arbitrary $g\colon a \to c$. Taking $K = \mathrm{pt}$, it says that $A \otimes f$ is in $\Sat$ whenever $f$ is and the tensors exist. Transfinite composites of $1$-morphisms in $\Sat$ are in $\Sat$ by (S2) and (S4).
Left cancellation also gives the ``reflection'' axiom of \cite[Prop.~1.17(4)]{dilibertilobbiasousa2022kz} without its second square: if $h$ is in $\Sat$, $l_{1}$, $l_{2}$ are coreflective and $hl_{1} \simeq l_{2}s$, then $s$ is in $\Sat$.

\begin{rmk}
\label{rmk:plain-vs-refl}
\gutentag{0079}%
None of this holds for $\overline{P}^{\dashv}$ (\cref{ex:saturation-fails}(3)). If every $p \in P$ and $u$ have cylinders, the plain saturation is recovered from the coreflective one: $u \in \overline{P}^{\dashv}$ if and only if $j_{u} \in \overline{J_{P}}$. Indeed, $\Pc^{\dashv}_{\emptyset,P} = \Pc_{\emptyset,J_{P}}$ by \cref{prop:cylinder}, and by \cref{prop:cylinder}(1) an object or $1$-morphism of $\D$ is $(\emptyset,\{u\})$-local in the plain sense if and only if it is $(\emptyset,\{j_{u}\})$-local.
\end{rmk}

\subsubsection*{Localizations and saturated classes}

\begin{prop}[Galois correspondence]
\label{prop:galois}
\gutentag{007A}%
Let $\C \subset \D$ be a locally full sub-$(\infty,2)$-category and $I$, $P$ classes of $1$-morphisms. Then $\C \subset \Pc_{I,P}$ if and only if $I$ is contained in $\Sat^{\mathrm{refl}}_{\C}$ and $P$ in $\Sat^{\mathrm{corefl}}_{\C}$. Consequently:
\begin{enumerate}
\item $\C \subset \Pc_{\Sat^{\mathrm{refl}}_{\C},\Sat^{\mathrm{corefl}}_{\C}}$, with equality if and only if $\C = \Pc_{I,P}$ for some $I$, $P$ (here $\Pc$ only depends on the $1$-morphisms in the saturations).
\item $\C \mapsto (\Sat^{\mathrm{refl}}_{\C},\Sat^{\mathrm{corefl}}_{\C})$ is injective on locally full sub-$(\infty,2)$-categories of the form $\Pc_{I,P}$, with left inverse $(\Sat,\Sat') \mapsto \Pc_{\Sat,\Sat'}$, and its values are pairs of a reflectively and a coreflectively saturated class (\cref{thm:saturation-saturated}).
\end{enumerate}
\end{prop}

\begin{proof}
\gutentag{007B}%
$\Pc_{I,P}$ is the intersection of the $\Pc_{\{i\},\emptyset}$, $i \in I$, and the $\Pc_{\emptyset,\{p\}}$, $p \in P$, so the first claim is \cref{lem:saturation-explicit}. For (1), apply it to $I = \Sat^{\mathrm{refl}}_{\C}$, $P = \Sat^{\mathrm{corefl}}_{\C}$; if $\C = \Pc_{I,P}$, then $I \subset \Sat^{\mathrm{refl}}_{\C}$ and $P \subset \Sat^{\mathrm{corefl}}_{\C}$ give the reverse inclusion. (2) follows from (1).
\end{proof}

\gutentag{007C}%
So $2$-dimensional localizations are determined by their saturated classes. The saturation of $(I,P)$ is a pair, and the two members depend on both $I$ and $P$: in general the $1$-morphisms in $\Sat^{\mathrm{corefl}}_{\Pc_{I,P}}$ strictly contain $\overline{P}$. By \cref{prop:2dloc-commute} the mixed case reduces to two pure ones in stages, provided the left adjoints exist, so we now fix $I = \emptyset$.
What \cref{prop:galois} does not say is which coreflectively saturated classes are saturations. In the one-dimensional case, for presentable $\C$ and a set $S$, this is the statement that the $S$-local equivalences form the strongly saturated class generated by $S$. Its two-dimensional analogue reduces to the units of the localization:

\begin{thm}
\label{thm:sat-units}
\gutentag{007D}%
Let $P$ be a class of $1$-morphisms such that $\iota\colon \C := \Pc_{\emptyset,P} \to \D$ has a left adjoint $L$ with unit $\eta\colon \mathrm{id}_{\D} \Rightarrow \iota L$ (by \cref{thm:2dloc-local}, which we will prove in \cref{subsec:2dloc-presentable}, this holds if $\D$ is presentable and $P$ is a set).
\begin{enumerate}
\item Every coreflectively saturated $\Sat$ containing all $\eta_{d}$, $d \in \D$, contains every $f \in \overline{P}$.
\item If every $\eta_{d}$ is in $\langle P \rangle$, then the $1$-morphisms in $\langle P \rangle$ form the class $\overline{P}$.
\item If every $\eta_{d}$ is in $\overline{P}$, then the $1$-morphisms in $\langle \{\eta_{d}\}_{d \in \D} \rangle$ form the class $\overline{P}$, and $\C = \Pc_{\emptyset,\{\eta_{d}\}_{d \in \D}}$.
\end{enumerate}
Dually for $\Pc_{I,\emptyset}$, reflectively saturated classes, and $\overline{I}$.
\end{thm}

\begin{proof}
\gutentag{007E}%
(1) Let $f\colon x \to y$ be in $\overline{P}$. By \cref{prop:saturation-L}(2), $Lf$ is coreflective in $\C$, hence $\iota Lf$ is coreflective in $\D$, as functors preserve adjunctions, and so $\iota Lf$ is in $\Sat$ by (S1). By naturality of $\eta$ and (S2), $\eta_{y}f \simeq \iota Lf \circ \eta_{x}$ is in $\Sat$. As $\eta_{y}$ is in $\Sat$, (S3) shows that $f$ is in $\Sat$.
(2) By \cref{thm:saturation-saturated}, $\Sat^{\mathrm{corefl}}_{\C}$ is coreflectively saturated and contains $P$, so the $1$-morphisms in $\langle P \rangle$ lie in $\overline{P}$ (\cref{lem:saturation-explicit}). The converse is (1) for $\Sat = \langle P \rangle$.
(3) As $\Sat^{\mathrm{corefl}}_{\C}$ is coreflectively saturated and contains every $\eta_{d}$, the $1$-morphisms in $\langle \{\eta_{d}\} \rangle$ lie in $\overline{P}$; the converse is (1). Let $\C' := \Pc_{\emptyset,\{\eta_{d}\}}$. As every $\eta_{d}$ is in $\Sat^{\mathrm{corefl}}_{\C}$, \cref{prop:galois} gives $\C \subset \C'$. Conversely, $\Sat^{\mathrm{corefl}}_{\C'}$ is coreflectively saturated and contains every $\eta_{d}$ (\cref{prop:galois}), so it contains $\langle \{\eta_{d}\} \rangle$ and hence $P \subset \overline{P}$; by \cref{prop:galois}, $\C' \subset \Pc_{\emptyset,P} = \C$.
\end{proof}

\begin{rmk}
\label{rmk:sat-units-kz}
\gutentag{007F}%
By \cref{prop:saturation-L}(2), the hypothesis of (3) says that every $L\eta_{d}$ is coreflective in $\C$. This is a form of lax idempotence of the monad $\iota L$; e.g.\ for $\D = \Cat$ and $P = \{\emptyset \to \mathrm{pt}\}$, where $\iota L$ freely adjoins an initial object, the right adjoint of $L\eta_{d}\colon d^{\triangleleft} \to (d^{\triangleleft})^{\triangleleft}$ collapses the two initial objects.
The hypothesis of (2) implies that of (3), and it would follow from a two-dimensional small object argument producing $\eta_{d}$ from $P$ by the operations (S1)--(S4).
For $2$-categories, Di Liberti--Lobbia--Sousa \cite[Thm.~4.3]{dilibertilobbiasousa2022kz} construct $\eta_{d}$ as a transfinite composite of wide bipushouts of bipushouts along $1$-morphisms of $P$ and of bicoequifiers, and show that every $\eta_{d}$ lies in $\overline{P}$ \cite[Lemma~3.5]{dilibertilobbiasousa2022kz}, which gives the hypothesis of (3). Pushouts, wide pushouts and transfinite composites are instances of (S2) and (S4), but they verify the bicoequifier and limit steps by hand, so their argument does not give the hypothesis of (2).
\end{rmk}

\begin{rmk}
\label{rmk:sat-open}
\gutentag{0080}%
What remains open for the correspondence:
\begin{enumerate}
\item whether every $\eta_{d}$ lies in $\overline{P}$ for $(\infty,2)$-categories (lax idempotence, \cref{rmk:sat-units-kz});
\item whether it lies in $\langle P \rangle$, which by \cref{thm:sat-units}(2) would identify $\overline{P}$ with the saturated class generated by $P$ on $1$-morphisms. The ingredient missing is an $(\infty,2)$-categorical replacement for the bicoequifier steps, which in an $(\infty,2)$-category have to control all higher morphisms of the hom-categories;
\item the analogous statement for squares. \Cref{thm:sat-units} compares only $1$-morphisms. Its proof uses left cancellation for $1$-morphisms, and right Beck--Chevalley squares do not satisfy the analogous cancellation: in \cref{ex:saturation-fails}(2), $\tau$ factors as $\sigma$ followed by the square $\mathrm{id}_{\Delta^{1}} \to \mathrm{id}_{\mathrm{pt}}$ with both components $\Delta^{1} \to \mathrm{pt}$, and $\tau$ and the latter are right Beck--Chevalley while $\sigma$ is not. So we do not know whether $\langle P \rangle$ contains all $\Sat^{\mathrm{corefl}}_{\C}$-squares.
\end{enumerate}
\end{rmk}
